\chapter{La \texorpdfstring{$\beta$}{Beta}-élimination}\label{ch:b1:elimination}

Traitons du 2-bromo-2-méthylpropane par l'éthanolate de sodium dans
l'éthanol à température ambiante~: le produit principal est un éther~;
chauffons le même mélange~: le produit principal est un alcène, le
2-méthylpropène, qui ne contient aucun groupe éthoxy. L'ion éthanolate a
agi non comme un \omterm{def:b1:electronic-effects:nucleophile}{nucléophile}, attaquant un carbone, mais comme une base,
arrachant un proton au carbone voisin pendant que le bromure partait.
Cette compétition entre substitution et élimination traverse toute la
synthèse organique. Ce chapitre décrit les deux mécanismes
d'élimination, leur remarquable exigence géométrique, la règle qui
prévoit quel alcène se forme, et la façon d'orienter une réaction dans
un sens ou dans l'autre.

\begin{recall}
SN1 et SN2 (\cref{ch:b1:nucleophilic-substitution})~; \omterm{def:b1:stereochemistry-in-depth:representations}{projections de
Newman}, \omterm{def:b1:stereochemistry-in-depth:isomers}{conformations} du cyclohexane et descripteurs $Z$/$E$
(\cref{ch:b1:stereochemistry-in-depth})~; \omterm{def:b1:electronic-effects:intermediates}{carbocations} et bases
(\cref{ch:b1:electronic-effects}).
\end{recall}

\section{Le mécanisme E2}

\begin{definition}[\texorpdfstring{$\beta$}{Beta}-élimination, E2 et E1]\label{def:b1:elimination:elimination}
Une \emph{$\beta$-élimination}\index{beta-elimination@$\beta$-élimination} retire un \omterm{def:b1:electronic-effects:nucleophile}{groupe partant} Y
d'un carbone (le carbone $\alpha$) et un proton d'un carbone voisin (un
carbone $\beta$), en formant une liaison double entre eux. Dans le
\emph{mécanisme E2}\index{mecanisme E2@mécanisme E2}, une base arrache le proton pendant que Y
s'en va, en un seul \omterm{def:b1:elementary-steps:elementary-step}{acte élémentaire}~; dans le \emph{mécanisme
E1}\index{mecanisme E1@mécanisme E1}, Y part d'abord, ce qui donne un \omterm{def:b1:electronic-effects:intermediates}{carbocation}, qui perd
ensuite un proton en $\beta$.
\end{definition}

\begin{proposition}[Loi de vitesse de l'E2]\label{prop:b1:elimination:e2-law}
Pour l'E2, $v = k[\mathrm{RY}][\mathrm{B}]$, où B est la base.
\end{proposition}

\begin{proof}
Un seul \omterm{def:b1:elementary-steps:elementary-step}{acte élémentaire} bimoléculaire, faisant intervenir le \omterm{def:b1:nucleophilic-substitution:substitution}{substrat}
et la base~: c'est la loi d'action des masses des \omterm{def:b1:elementary-steps:elementary-step}{actes élémentaires}
(\cref{ch:b1:elementary-steps}).
\end{proof}

\begin{definition}[Antipériplanaire et synpériplanaire]\label{def:b1:elimination:periplanar}
Deux liaisons portées par des atomes voisins sont
\emph{antipériplanaires}\index{antiperiplanaire@antipériplanaire} lorsqu'elles sont dans un même plan, de
part et d'autre (\omterm{def:b1:stereochemistry-in-depth:conformations}{angle de torsion} $180^\circ$), et
\emph{synpériplanaires}\index{synperiplanaire@synpériplanaire} lorsqu'elles sont dans un même plan, du
même côté ($0^\circ$).
\end{definition}

\begin{omfigure}
\begin{tikzpicture}[font=\small, >={Stealth}]
  \pic (n) at (0,0) {omnewman={90}{270}};
  \node at (n-f1) {H}; \node at (n-f2) {\ce{CH3}}; \node at (n-f3) {H};
  \node at (n-b1) {Br}; \node at (n-b2) {\ce{CH3}}; \node at (n-b3) {H};
  \node (b) at (-2.6,2.0) {\ce{EtO-}};
  \draw[->, thick, omThm] (b) to[bend left=20] ($(n-f1)+(-0.15,0.1)$);
  \draw[->, thick, omThm] (0.12,0.8) to[bend left=40] (0.12,0.25);
  \draw[->, thick, omThm] (0.12,-0.6) to[bend right=35] ($(n-b1)+(0.25,0.15)$);
  \node[below] at (0,-2.0) {H et Br antipériplanaires};
  \draw[->, thick] (2.4,0) -- (3.6,0) node[midway, above] {E2};
  \node at (5.6,0) {\chemfig{C(-[:150]H_3C)(-[:210]H)=C(-[:30]H)(-[:-30]CH_3)}};
  \node at (5.6,-0.9) {\cip{E}-but-2-ène};
\end{tikzpicture}

{\small Réaction E2 du 2-bromobutane, vue le long de la liaison C3--C2
(C3 en avant, C2, qui porte Br, en arrière) dans l'une de ses deux
\omterm{def:b1:stereochemistry-in-depth:isomers}{conformations} \omterm{def:b1:elimination:periplanar}{antipériplanaires}~; l'autre, où l'autre hydrogène de C3
est anti par rapport à Br, donne le \cip{Z}-but-2-ène. Trois doublets
d'électrons se déplacent ensemble~: la base arrache le proton, le
doublet de C--H devient la liaison $\pi$, le doublet de C--Br part avec
le brome. Les groupes qui se font face dans la \omterm{def:b1:stereochemistry-in-depth:representations}{projection de Newman} se
retrouvent du même côté de la liaison double.}
\end{omfigure}

\begin{proposition}[L'E2 est anti et stéréospécifique]\label{prop:b1:elimination:anti}
L'E2 se fait à partir de la \omterm{def:b1:stereochemistry-in-depth:isomers}{conformation} dans laquelle les liaisons
C--H et C--Y sont \omterm{def:b1:elimination:periplanar}{antipériplanaires}. Elle est stéréospécifique~: des
\omterm{def:b1:nucleophilic-substitution:substitution}{substrats} \omterm{def:b1:stereochemistry-in-depth:enantiomers}{diastéréoisomères} donnent des \omterm{def:b1:stereochemistry-in-depth:isomers}{stéréoisomères} différents de
l'alcène, fixés par la disposition anti.
\end{proposition}

\begin{proof}
Dans l'\omterm{def:b1:elementary-steps:energy-profile}{état de transition}, le \omterm{def:b1:lewis-resonance-vsepr:lewis}{doublet liant} de C--H doit pénétrer dans
la région antiliante de la liaison C--Y pendant que les deux carbones
deviennent trigonaux~; le recouvrement est le meilleur lorsque les deux
liaisons sont parallèles, dans un même plan. La disposition anti est
décalée (de basse énergie) et maintient la base loin du \omterm{def:b1:electronic-effects:nucleophile}{groupe
partant}~; la disposition syn est éclipsée. H et Y étant anti, les deux
autres groupes de chaque carbone sont fixés~: ceux qui sont du même côté
dans la \omterm{def:b1:stereochemistry-in-depth:representations}{projection de Newman} se retrouvent \textit{cis} sur la liaison
double. Un diastéréoisomère a deux groupes échangés sur un carbone, et
donne l'autre alcène.
\end{proof}

\begin{proposition}[E2 sur un cyclohexane]\label{prop:b1:elimination:cyclohexane-e2}
Sur un cycle cyclohexane, l'E2 exige que le \omterm{def:b1:electronic-effects:nucleophile}{groupe partant} et
l'hydrogène en $\beta$ soient tous deux \omterm{def:b1:stereochemistry-in-depth:conformations}{axiaux} (\textit{trans}-diaxiaux)~;
un \omterm{def:b1:electronic-effects:nucleophile}{groupe partant} maintenu \omterm{def:b1:stereochemistry-in-depth:conformations}{équatorial} ne peut pas réagir tant que le
cycle ne s'inverse pas.
\end{proposition}

\begin{proof}
Sur des carbones adjacents d'une \omterm{def:b1:stereochemistry-in-depth:conformations}{chaise}, deux liaisons \omterm{def:b1:stereochemistry-in-depth:conformations}{axiales} sont
\omterm{def:b1:elimination:periplanar}{antipériplanaires} (l'une vers le haut, l'autre vers le bas, parallèles à
l'axe)~; une liaison \omterm{def:b1:stereochemistry-in-depth:conformations}{équatoriale} n'est anti que par rapport à des
liaisons C--C du cycle, jamais par rapport à une liaison C--H.
\end{proof}

\begin{omfigure}
\begin{tikzpicture}[font=\small]
  \pic (a) at (0,0) {omchair};
  \draw[thick, omThm] (a-c1) -- (a-a1) node[above] {Cl};
  \draw[thick] (a-c1) -- (a-e1) node[right] {H};
  \draw[thick, omDef] (a-c2) -- (a-a2) node[below] {H};
  \draw[thick, omDef] (a-c6) -- (a-a6) node[below] {H};
  \node[align=left, anchor=west] at (3.0,0.2) {Cl axial (vers le haut) sur C1~;\\H axial (vers le bas) sur les deux voisins~:\\deux H antipériplanaires};
\end{tikzpicture}

{\small Disposition \textit{trans}-diaxiale sur une \omterm{def:b1:stereochemistry-in-depth:conformations}{chaise}~: le chlore
\omterm{def:b1:stereochemistry-in-depth:conformations}{axial} est \omterm{def:b1:elimination:periplanar}{antipériplanaire} aux hydrogènes \omterm{def:b1:stereochemistry-in-depth:conformations}{axiaux} des deux carbones
voisins (en bleu).}
\end{omfigure}

\section{Le mécanisme E1}

\begin{proposition}[Loi de vitesse de l'E1]\label{prop:b1:elimination:e1-law}
Pour l'E1, $v = k[\mathrm{RY}]$, la vitesse étant celle de l'ionisation.
\end{proposition}

\begin{proof}
Comme pour la SN1 (\cref{prop:b1:nucleophilic-substitution:sn1-law})~:
le \omterm{def:b1:electronic-effects:intermediates}{carbocation} est formé au cours d'une étape lente et consommé
rapidement, ici par perte d'un proton au profit du solvant ou d'une \omterm{def:b1:predominant-reaction:strong-acid}{base
faible}~; l'\omterm{def:b1:elementary-steps:qssa}{approximation des états quasi stationnaires} laisse
$v = k_1[\mathrm{RY}]$.
\end{proof}

L'E1 partage son \omterm{def:b1:electronic-effects:intermediates}{carbocation} avec la SN1~: les \omterm{def:b1:nucleophilic-substitution:substitution}{substrats} tertiaires,
dans les \omterm{def:b1:intermolecular-forces-solvents:solvent-classes}{solvants protiques}, donnent les deux, l'alcène étant favorisé
par le chauffage. Le \omterm{def:b1:electronic-effects:intermediates}{carbocation} étant plan et le proton n'étant perdu
qu'ensuite, l'E1 n'a pas d'exigence géométrique et n'est pas
stéréospécifique~; elle donne surtout l'alcène le plus stable, en
général l'alcène \cip{E}.

\begin{omfigure}
\setchemfig{atom sep=2.1em}
\schemestart
\chemfig{C(-[2]CH_3)(-[4]H_3C)(-[6]CH_3)-Br}
\arrow{->[lente]}
\chemfig{C^{+}(-[2]CH_3)(-[:210]H_3C)(-[:-30]CH_3)}
\arrow{->[rapide][\ce{-H+}]}
\chemfig{C(-[2]CH_3)(-[:210]H_3C)=[:-30]CH_2}
\schemestop
\par\vspace{4pt}

{\small E1~: ionisation du 2-bromo-2-méthylpropane en \omterm{def:b1:electronic-effects:intermediates}{carbocation}, puis
perte d'un proton d'un groupe méthyle au profit du solvant, ce qui donne
le 2-méthylpropène.}
\end{omfigure}

\section{Régiosélectivité~: la règle de Zaïtsev}

\begin{definition}[Réactions régiosélectives et stéréosélectives]\label{def:b1:elimination:selectivity}
Une réaction qui peut donner plusieurs \omterm{def:b1:stereochemistry-in-depth:isomers}{isomères de constitution} mais en
forme un majoritairement est une \emph{réaction
régiosélective}\index{reaction regioselective@réaction régiosélective}~; une réaction qui peut donner plusieurs
\omterm{def:b1:stereochemistry-in-depth:isomers}{stéréoisomères} mais en forme un majoritairement est une \emph{réaction
stéréosélective}\index{reaction stereoselective@réaction stéréosélective}.
\end{definition}

\begin{proposition}[Règle de Zaïtsev]\label{prop:b1:elimination:zaitsev}
Lorsque plusieurs hydrogènes en $\beta$ peuvent être arrachés, l'alcène
majoritaire est en général le plus substitué (le plus stable), et
l'isomère \cip{E} plutôt que l'isomère \cip{Z} (\emph{règle de
Zaïtsev}\index{regle de Zaitsev@règle de Zaïtsev}). Les bases encombrées et les contraintes
géométriques, comme l'exigence \textit{trans}-diaxiale, peuvent
l'emporter sur elle.
\end{proposition}

\begin{proof}
Les groupes alkyle portés par les carbones d'une liaison double la
stabilisent (comme ils stabilisent les \omterm{def:b1:electronic-effects:intermediates}{carbocations}), et les alcènes
\cip{E} évitent l'encombrement des groupes \textit{cis}. Les \omterm{def:b1:elementary-steps:energy-profile}{états de
transition} de l'E2 et de l'E1 ont un caractère partiel de liaison
double, si bien que l'alcène le plus stable se forme le plus vite. Une
base encombrée atteint plus facilement les hydrogènes les moins
encombrés~; et un hydrogène qui ne peut pas être \omterm{def:b1:elimination:periplanar}{antipériplanaire} au
\omterm{def:b1:electronic-effects:nucleophile}{groupe partant} n'est pas du tout arraché par E2, quelle que soit la
stabilité de l'alcène qu'il donnerait.
\end{proof}

\begin{method}[Prévoir l'alcène]\label{met:b1:elimination:predict-alkene}
\begin{enumerate}
  \item Recenser les carbones $\beta$ qui portent des hydrogènes.
  \item Pour l'E2, ne garder que les hydrogènes qui peuvent être
        \omterm{def:b1:elimination:periplanar}{antipériplanaires} à Y (sur un cycle~: \textit{trans}-diaxiaux,
        dans une \omterm{def:b1:stereochemistry-in-depth:conformations}{chaise} qui peut réellement se former).
  \item Parmi eux, privilégier celui qui donne l'alcène le plus
        substitué (Zaïtsev), sauf si la base est encombrée.
  \item Pour chacun, dessiner la \omterm{def:b1:stereochemistry-in-depth:representations}{projection de Newman} en \omterm{def:b1:stereochemistry-in-depth:conformations}{conformation
        anti} pour lire la géométrie \cip{E}/\cip{Z}.
\end{enumerate}
\end{method}

\section{Substitution ou élimination~?}

\begin{proposition}[Substitution contre élimination]\label{prop:b1:elimination:sn-vs-e}
L'élimination est favorisée par les \omterm{def:b1:predominant-reaction:strong-acid}{bases fortes} et encombrées, par une
température élevée et par les \omterm{def:b1:nucleophilic-substitution:substitution}{substrats} tertiaires ou encombrés~; la
substitution, par les bons \omterm{def:b1:electronic-effects:nucleophile}{nucléophiles} qui sont des \omterm{def:b1:predominant-reaction:strong-acid}{bases faibles}, par
une température basse et par les \omterm{def:b1:nucleophilic-substitution:substitution}{substrats} primaires.
\end{proposition}

\begin{proof}
La force de la base favorise l'attaque sur l'hydrogène~; l'encombrement
gêne l'attaque sur un carbone encombré, mais pas sur un hydrogène
exposé. Un \omterm{def:b1:electronic-effects:carbon-class}{carbone tertiaire} est inaccessible à la SN2, et son
\omterm{def:b1:electronic-effects:intermediates}{carbocation} perd facilement des protons. L'élimination transforme une
molécule en deux ou trois (alcène, \omterm{def:b1:electronic-effects:nucleophile}{groupe partant}, base protonée)~: elle
est favorisée quand la température augmente, résultat justifié dans le
volume de deuxième année.
\end{proof}

\begin{omfigure}
\begin{tabular}{llll}
\hline
\omterm{def:b1:nucleophilic-substitution:substitution}{substrat} & \omterm{def:b1:predominant-reaction:strong-acid}{base forte}, petite & \omterm{def:b1:predominant-reaction:strong-acid}{base forte}, encombrée & \omterm{def:b1:predominant-reaction:strong-acid}{base faible}, bon \omterm{def:b1:electronic-effects:nucleophile}{nucléophile} \\
\hline
primaire & SN2 (un peu d'E2) & E2 & SN2 \\
secondaire & E2 (et SN2) & E2 & SN2 (aprotique), SN1/E1 (protique) \\
tertiaire & E2 & E2 & SN1/E1 (protique) \\
\hline
\end{tabular}

\smallskip
{\small Une première orientation~: la réaction principale dans chaque
cas. Le chauffage favorise l'élimination, colonne par colonne.}
\end{omfigure}

\section{Exercices}

\begin{exercise}[$\star$]\label{exo:b1:elimination:1}
Donner les alcènes qui peuvent se former par élimination de HBr à
partir du 2-bromobutane, du 2-bromo-2-méthylbutane et du
bromocyclohexane, ainsi que celui que la \omterm{prop:b1:elimination:zaitsev}{règle de Zaïtsev} prévoit
majoritaire.
\end{exercise}

\begin{exercise}[$\star$]\label{exo:b1:elimination:2}
Écrire les \omterm{def:b1:rate-laws:order}{lois de vitesse} de l'E1 et de l'E2. Comment les
distinguer expérimentalement pour le 2-bromo-2-méthylpropane dans
l'éthanol contenant de l'éthanolate~?
\end{exercise}

\begin{exercise}[$\star$]\label{exo:b1:elimination:3}
Dessiner les \omterm{def:b1:stereochemistry-in-depth:representations}{projections de Newman} du 2-bromobutane le long de C2--C3
dans ses trois \omterm{def:b1:stereochemistry-in-depth:conformations}{conformations décalées}, et dire lesquelles peuvent subir
une E2.
\end{exercise}

\begin{exercise}[$\star$]\label{exo:b1:elimination:4}
Prévoir la réaction principale (SN1, SN2, E1, E2)~: bromoéthane avec
l'hydroxyde de sodium dans l'eau à \qty{25}{\celsius}~;
2-bromo-2-méthylpropane avec le \textit{tert}-butanolate de potassium~;
2-bromopropane avec l'iodure de sodium dans la propanone~;
2-bromo-2-méthylpropane dans l'éthanol chaud.
\end{exercise}

\begin{exercise}[$\star\star$]\label{exo:b1:elimination:5}
Montrer, à l'aide de \omterm{def:b1:stereochemistry-in-depth:representations}{projections de Newman}, que l'élimination E2 de HBr
à partir des deux \omterm{def:b1:stereochemistry-in-depth:enantiomers}{diastéréoisomères} du 2-bromo-3-méthylpentane qui
peuvent perdre l'hydrogène de C3 donne les deux \omterm{def:b1:stereochemistry-in-depth:isomers}{stéréoisomères} du
3-méthylpent-2-ène.
\end{exercise}

\begin{exercise}[$\star\star$]\label{exo:b1:elimination:6}
Écrire le mécanisme de la déshydratation du 2-méthylbutan-2-ol dans
l'acide sulfurique concentré chaud (E1) et donner l'alcène majoritaire.
\end{exercise}

\begin{exercise}[$\star\star$]\label{exo:b1:elimination:7}
Expliquer pourquoi le \textit{trans}-1-bromo-4-\textit{tert}-butylcyclohexane
réagit très lentement par E2 alors que son isomère \textit{cis} réagit
rapidement. (Le groupe \textit{tert}-butyle reste \omterm{def:b1:stereochemistry-in-depth:conformations}{équatorial}.)
\end{exercise}

\begin{exercise}[$\star\star$]\label{exo:b1:elimination:8}
Le 2-bromo-2-méthylbutane donne, avec l'éthanolate de sodium,
surtout du 2-méthylbut-2-ène~; avec le \textit{tert}-butanolate de
potassium, surtout du 2-méthylbut-1-ène. Expliquer.
\end{exercise}

\begin{exercise}[$\star\star$]\label{exo:b1:elimination:9}
Pourquoi n'observe-t-on jamais d'élimination avec le bromométhane~? Et
avec le 1-bromo-2,2-diméthylpropane par E2~?
\end{exercise}

\begin{exercise}[$\star\star\star$]\label{exo:b1:elimination:10}
Dans l'éthanol à \qty{25}{\celsius}, le 2-bromo-2-méthylpropane donne un
mélange de l'éther (SN1) et de l'alcène (E1) dont les proportions ne
dépendent pas de la concentration du \omterm{def:b1:nucleophilic-substitution:substitution}{substrat}. Montrer que les deux
produits proviennent du même \omterm{def:b1:electronic-effects:intermediates}{carbocation} et que leur rapport est le
rapport de deux \omterm{def:b1:rate-laws:order}{constantes de vitesse}.
\end{exercise}

\begin{exercise}[$\star\star\star$]\label{exo:b1:elimination:11}
L'un des \omterm{def:b1:stereochemistry-in-depth:enantiomers}{diastéréoisomères} du 1,2-dibromo-1,2-diphényléthane (la forme
méso) donne, par E2 avec l'éthanolate, un seul stéréoisomère du
1-bromo-1,2-diphényléthène. Trouver lequel, à l'aide d'une \omterm{def:b1:stereochemistry-in-depth:representations}{projection de
Newman}.
\end{exercise}

\begin{exercise}[$\star\star\star$]\label{exo:b1:elimination:12}
À l'aide des énergies du \cref{ch:b1:stereochemistry-in-depth},
expliquer pourquoi une réaction E2 sur un cyclohexane peut être ralentie
d'un facteur de plusieurs centaines lorsque le \omterm{def:b1:electronic-effects:nucleophile}{groupe partant} doit
devenir \omterm{def:b1:stereochemistry-in-depth:conformations}{axial} dans une \omterm{def:b1:stereochemistry-in-depth:conformations}{chaise} encombrée. Exprimer la vitesse comme le
produit d'une fraction de conformère par une \omterm{def:b1:rate-laws:order}{constante de vitesse}.
\end{exercise}

\section{Problème~: les chlorures de menthyle et de néomenthyle}

\begin{problem}[{Problème du week-end --- les chaises de deux chlorures
diastéréoisomères, les hydrogènes antipériplanaires dont dispose chacun,
les alcènes formés, et le rapport de leurs vitesses d'E2}]
\label{pb:b1:elimination:1}
Le chlorure de menthyle et le chlorure de néomenthyle sont les chlorures
du menthol et du néomenthol (\cref{ch:b1:stereochemistry-in-depth})~:
sur le cycle, C1 porte Cl, C2 le groupe isopropyle, C5 le groupe méthyle.
Dans le chlorure de menthyle, les trois groupes peuvent être tous
\omterm{def:b1:stereochemistry-in-depth:conformations}{équatoriaux}~; dans le chlorure de néomenthyle, Cl est \omterm{def:b1:stereochemistry-in-depth:conformations}{axial} lorsque les
deux groupes alkyle sont \omterm{def:b1:stereochemistry-in-depth:conformations}{équatoriaux}. On les traite tous deux par
l'éthanolate de sodium dans l'éthanol, une réaction E2. Données du
problème~: \qty{95}{\percent} du chlorure de néomenthyle est dans la
\omterm{def:b1:stereochemistry-in-depth:conformations}{chaise} où Cl est \omterm{def:b1:stereochemistry-in-depth:conformations}{axial}~; pour le chlorure de menthyle, la \omterm{def:b1:stereochemistry-in-depth:conformations}{chaise} où Cl
est \omterm{def:b1:stereochemistry-in-depth:conformations}{axial} (les trois groupes \omterm{def:b1:stereochemistry-in-depth:conformations}{axiaux}) représente une fraction
\num{5.0e-3}~; on admet que chaque hydrogène \omterm{def:b1:elimination:periplanar}{antipériplanaire} d'une
\omterm{def:b1:stereochemistry-in-depth:conformations}{chaise} réactive réagit avec la même \omterm{def:b1:rate-laws:order}{constante de vitesse}. Le chlorure de
néomenthyle donne \qty{78}{\percent} d'alcène trisubstitué.

\medskip
\noindent\textbf{Partie I --- Les \omterm{def:b1:stereochemistry-in-depth:conformations}{chaises}.}
\begin{enumerate}
  \item Dessiner la \omterm{def:b1:stereochemistry-in-depth:conformations}{chaise} préférée du chlorure de menthyle.
  \item Dessiner la \omterm{def:b1:stereochemistry-in-depth:conformations}{chaise} préférée du chlorure de néomenthyle.
  \item Les deux chlorures sont-ils \omterm{def:b1:stereochemistry-in-depth:enantiomers}{énantiomères} ou \omterm{def:b1:stereochemistry-in-depth:enantiomers}{diastéréoisomères}~?
  \item Énoncer l'exigence géométrique de l'E2 sur un cycle.
  \item Dans quelle \omterm{def:b1:stereochemistry-in-depth:conformations}{chaise} le chlorure de menthyle peut-il réagir~? La
        dessiner.
  \item Pourquoi cette \omterm{def:b1:stereochemistry-in-depth:conformations}{chaise} est-elle rare~? Utiliser le coût
        énergétique des groupes \omterm{def:b1:stereochemistry-in-depth:conformations}{axiaux}.
\end{enumerate}

\medskip
\noindent\textbf{Partie II --- Les hydrogènes anti.}
\begin{enumerate}[resume]
  \item Recenser les carbones $\beta$ du chlorure (C2 et C6) et leurs
        hydrogènes.
  \item Dans la \omterm{def:b1:stereochemistry-in-depth:conformations}{chaise} réactive du chlorure de néomenthyle, quels
        hydrogènes sont \omterm{def:b1:elimination:periplanar}{antipériplanaires} à Cl~?
  \item Dans la \omterm{def:b1:stereochemistry-in-depth:conformations}{chaise} réactive du chlorure de menthyle, l'hydrogène de
        C2 est-il \omterm{def:b1:stereochemistry-in-depth:conformations}{axial}~? Un hydrogène de C6 est-il \omterm{def:b1:stereochemistry-in-depth:conformations}{axial}~?
  \item Combien d'hydrogènes \omterm{def:b1:elimination:periplanar}{antipériplanaires} chaque chlorure
        offre-t-il~?
  \item Dessiner une \omterm{def:b1:stereochemistry-in-depth:representations}{projection de Newman} le long de C1--C6 pour la
        \omterm{def:b1:stereochemistry-in-depth:conformations}{chaise} réactive du chlorure de menthyle.
  \item Pourquoi une élimination \omterm{def:b1:elimination:periplanar}{synpériplanaire} ne pourrait-elle pas
        sauver la réaction~?
\end{enumerate}

\medskip
\noindent\textbf{Partie III --- Les produits.}
\begin{enumerate}[resume]
  \item Nommer l'alcène formé par arrachement de l'hydrogène de C2 et
        celui formé par arrachement d'un hydrogène de C6.
  \item Lequel la \omterm{prop:b1:elimination:zaitsev}{règle de Zaïtsev} favorise-t-elle, et pourquoi~?
  \item Que donne le chlorure de néomenthyle~? Le rapport 78/22
        est-il cohérent avec la \omterm{prop:b1:elimination:zaitsev}{règle de Zaïtsev}~?
  \item Que donne le chlorure de menthyle~? Est-ce cohérent avec la
        \omterm{prop:b1:elimination:zaitsev}{règle de Zaïtsev}~?
  \item Quelle réaction est stéréospécifique, régiosélective, ou les
        deux~?
  \item Pourquoi la substitution (SN2) est-elle minoritaire ici avec
        l'éthanolate~?
\end{enumerate}

\medskip
\noindent\textbf{Partie IV --- Les vitesses.}
\begin{enumerate}[resume]
  \item Écrire la vitesse de chaque réaction en fonction de la fraction
        de \omterm{def:b1:stereochemistry-in-depth:conformations}{chaise} réactive et du nombre d'hydrogènes \omterm{def:b1:elimination:periplanar}{antipériplanaires}.
  \item Calculer la vitesse relative du chlorure de néomenthyle.
  \item Calculer la vitesse relative du chlorure de menthyle.
  \item Calculer le rapport des vitesses $k(\text{néomenthyle})/k(\text{menthyle})$.
\end{enumerate}
\end{problem}
